\documentclass[10pt]{article}
\usepackage[margin=2cm]{geometry}
\usepackage{booktabs}
\usepackage{amsmath}
\usepackage{enumitem}
\usepackage[most]{tcolorbox}
\usepackage{hyperref}
\setlength{\emergencystretch}{3em}
\setlength{\parindent}{0pt}
\setlength{\parskip}{5pt}

% Grey box for the rubric of each part
\newenvironment{rubric}{%
  \begin{tcolorbox}[colback=black!4, colframe=black!40, boxrule=0.4pt,
                    arc=0pt, title={Rubric}, fonttitle=\bfseries\small,
                    coltitle=black, colbacktitle=black!12]
  \small\begin{itemize}[leftmargin=*, nosep]}{%
  \end{itemize}\end{tcolorbox}}

\begin{document}

\begin{center}
{\Large\bfseries CSCI 5942: Mock Quiz 6, Answers and Rubric}\\[3pt]
{\large Assignment 6: Serving GPT-2}
\end{center}

Each part is worth 0.75 points. The rubric follows the real quizzes:
\textbf{mechanism + numbers + a concrete verdict}. A correct verdict
without the numbers, or numbers without the reason, gets partial credit.

% =====================================================================
\section*{Set A}

\subsection*{1(a) Sizing the cache}

\begin{align*}
\text{KV per token} &= 2\ (\text{K, V}) \times 36 \times 1280 \times 2\,\text{B}
  = \mathbf{184{,}320\,B}\ (180\,\text{KiB}) \\
\text{contiguous, per request} &= 184{,}320 \times 1{,}024 = 188.7\,\text{MB}
  \;\Rightarrow\; 20 \times 10^9 / 188.7 \times 10^6 = 106.0
  \;\Rightarrow\; \mathbf{105}\ \text{requests} \\
\text{one block} &= 16 \times 184{,}320 = 2{,}949{,}120\,\text{B}
  \;\Rightarrow\; 20 \times 10^9 / 2{,}949{,}120 = \mathbf{6{,}781}\ \text{blocks} \\
\text{paged, per request} &= \lceil (70 + 40) / 16 \rceil = 7\ \text{blocks}
  \;\Rightarrow\; 6{,}781 / 7 = \mathbf{968}\ \text{requests}\ (\approx 9.2\times)
\end{align*}

\begin{rubric}
\item 0.25: KV bytes per token with the formula (2 $\times$ layers $\times$ width $\times$ bytes)
\item 0.25: contiguous count from 1{,}024 reserved slots
\item 0.25: paged count from the request's real length, rounded up to whole blocks
\end{rubric}

\textbf{Common mistakes.} Forgetting the factor 2 for K and V; using
4 bytes (fp32) instead of 2 (bf16); reserving only 110 slots for the
contiguous cache (it must reserve the whole context, because it does not
know the output length in advance).

\subsection*{1(b) The SLO decides \texttt{max\_batch}}

The largest batch whose step fits in 30\,ms is \textbf{$B = 256$}
(24.5\,ms; $B = 512$ takes 41\,ms). That gives $256 / 0.0245\,\text{s}
\approx$ \textbf{10{,}400 tokens/s}. The contiguous cache can only run 105
requests: between $B = 64$ and 128, the step is about 16.3\,ms, so about
$105 / 0.0163 \approx$ \textbf{6{,}400 tokens/s}. Paging buys about
\textbf{$1.6\times$}, not $9\times$.

Response to product: paging removes the \textbf{memory} limit, but
throughput stops growing once another limit binds. Up to about $B = 128$,
decode is memory-bound (the step time barely moves, 14.0 to 17.0\,ms), so
a bigger batch is almost free. Past that, the step becomes compute-bound
and the step time rises with the batch, and at $B = 256$ the \textbf{TBT
SLO} is the limit. The extra memory ($968 - 256$ requests' worth) can be
used for something else, such as prefix caching or longer contexts.

\begin{rubric}
\item 0.25: picks $B = 256$ from the table and the SLO
\item 0.25: tokens/s for both caches (about 10.4K vs.\ about 6.4K)
\item 0.25: explains why it is not $9\times$: memory-bound $\to$
  compute-bound, and the SLO now binds, not memory
\end{rubric}

\subsection*{2(a) Keep optimistic admission}

\textbf{Disagree.} Optimistic admission raises throughput from 2{,}780 to
3{,}350 tokens/s (+20\%) and the mean batch from 45.2 to 62.5 (+38\%), and
it lowers both mean latency (2.9 to 2.4\,s) and p90 (4.9 to 4.2\,s). The
safe engine reserves 7 blocks for every request's full length, so at
1\,GB (339 blocks) it can hold at most 48 requests, and most of those
reserved blocks are still empty. Optimistic admission reserves only the
next token, so more requests run per step.

A preemption costs less than it looks because recomputation is
\textbf{one prefill} of the prompt plus the tokens generated so far.
Prefill processes all those tokens in one forward pass and uses the GPU's
compute well. The original tokens took many memory-bound decode steps. The
work is not repeated token by token.

Keep it on, but watch the worst case: a request preempted several times
has a much worse latency than the mean shows, and with an even smaller
pool, preemptions could repeat (thrashing) and throughput could fall below
safe admission.

\begin{rubric}
\item 0.25: verdict with numbers (tok/s and mean batch both up, latency down)
\item 0.25: why recompute is cheap (one compute-efficient prefill vs.\ many
  memory-bound decode steps)
\item 0.25: names the risk (tail latency of preempted requests, or
  thrashing at smaller budgets)
\end{rubric}

\subsection*{2(b) The short-prompt bug}

The \textbf{mask} is most likely wrong: it does not hide each row's padded
slots. A short row's block table is padded on the right with
\textbf{block id 0}, and block 0 is a real block that another sequence may
be using. The gather copies it into the short row's K and V. The mask must
allow slot $s$ only when $s \le \texttt{positions[b, t]}$, the row's
\textbf{own} position. If the new hire builds the mask from the longest
row's positions and shares it with every row (for example
\texttt{positions.max(0)}), the short rows attend to another sequence's
keys and values.

It only shows up for shorter prompts because the longest row has no
padding, so every slot it sees is its own. With one sequence, or prompts
of the same length, no row is padded, so the bug is hidden.

\begin{rubric}
\item 0.25: names the mask (or the padded table entries) as the cause
\item 0.25: explains that padded entries point to block 0, which holds
  another sequence's data
\item 0.25: explains why the longest row and equal-length batches are not
  affected
\end{rubric}

\textbf{Other acceptable answer (partial credit).} A wrong \texttt{slots}
computation would also corrupt the cache, but it would usually break
single-sequence tests too, so the mask is the better diagnosis.

% =====================================================================
\newpage
\section*{Set B}

\subsection*{3(a) Prefix caching and the timestamp}

Automatic prefix caching hashes \textbf{full 16-token blocks}, and each
hash includes its parent's hash. The fixed instruction is 50 tokens, so 3
full blocks (\textbf{48 tokens}) are identical in every request. The 4th
block mixes the last 2 instruction tokens with the attribute list, which
differs per listing. So 48 of 70 prompt tokens (\textbf{69\%}) can be
reused.

The timestamp breaks this because a token's K and V depend on
\textbf{every earlier token and its position}. A different first token
changes the first block's hash, and every later hash includes it, so no
block matches and the hit rate drops to 0\%. Fix: put anything that
changes per request (timestamps, user IDs) at the \textbf{end} of the
prompt, or log it outside the prompt.

\begin{rubric}
\item 0.25: 48 tokens / 69\%, with the full-block reason
\item 0.25: why a different first token breaks every later block (K, V
  depend on the whole prefix; chained hash)
\item 0.25: the fix (move changing content to the end)
\end{rubric}

\subsection*{3(b) Chunked prefill}

\textbf{Chunked prefill with stall-free batching} (Sarathi-Serve). Each
step has a token budget $\tau$: first add every running decode (100
tokens), then fill the rest of the budget with a chunk of the summary's
prompt. The largest $\tau$ whose step fits in 30\,ms is \textbf{$\tau =
512$} (27\,ms; 768 tokens take 38\,ms).

Each step then has room for $512 - 100 = 412$ prefill tokens, so the
1{,}000-token prompt takes 3 steps ($412 + 412 + 176$). Today, the one big
step holds 1{,}100 tokens, slower than 50\,ms, so every autocomplete user
waits more than $3\times$ a normal 15\,ms step. With chunks, no step goes
over 27\,ms.

The summary request pays with a \textbf{higher TTFT}: its prefill now
spans 3 steps instead of one. The last step holds $176 + 100 = 276$
tokens, so it takes about 18\,ms, for about $27 + 27 + 18 = 72$\,ms in
total. Its later chunks also re-read the KV of the earlier chunks, and
each extra step repeats the fixed per-step cost.

\begin{rubric}
\item 0.25: names chunked prefill / stall-free batching and why it works
  (decode steps are memory-bound, so there is spare compute for prefill
  tokens)
\item 0.25: $\tau = 512$ from the table and the SLO, with the number of
  chunks
\item 0.25: the cost to the summary request (higher TTFT, KV re-reads,
  repeated per-step overhead)
\end{rubric}

\subsection*{4(a) GQA and capacity}

\begin{align*}
\text{KV per token} &= 2 \times 32 \times 8\ (\text{KV heads}) \times 128
  \times 2\,\text{B} = \mathbf{131{,}072\,B}\ (128\,\text{KiB}) \\
\text{one 8K conversation} &= 131{,}072 \times 8{,}192 = 1.07\,\text{GB} \\
\text{KV budget} &= 80 - 16 - 4 = 60\,\text{GB}
  \;\Rightarrow\; 60 / 1.07 = 55.9 \;\Rightarrow\; \mathbf{55}\ \text{conversations}
\end{align*}

With full multi-head attention (32 KV heads), each token needs
\textbf{$4\times$} as much, 512\,KiB, so one conversation needs 4.29\,GB,
and only $60 / 4.29 = 13.97$, or \textbf{13 conversations}, fit. GQA
shares one K/V head across 4 query heads, which is why Llama-3.1-8B uses
it.

\begin{rubric}
\item 0.25: KV per token with \texttt{n\_kv\_heads} (not \texttt{n\_heads})
\item 0.25: the budget (60\,GB) and 55 conversations
\item 0.25: the MHA case (13) and why GQA shrinks the cache
\end{rubric}

\subsection*{4(b) Not recomputing the history}

Any two of:
\begin{itemize}
\item \textbf{Prefix caching on the GPU} (vLLM automatic prefix caching or
  SGLang RadixAttention): keep the conversation's blocks after the turn
  ends and reuse them when the next turn arrives with the same prefix.
  \emph{Limit:} HBM is small, and the blocks are evicted (LRU) long before
  the user returns hours later.
\item \textbf{Hierarchical KV cache} (offload to CPU DRAM, SSD or a remote
  store, e.g.\ LMCache, Mooncake): much more capacity, so the history
  survives until the next turn. \emph{Limit:} loading costs PCIe or
  network bandwidth, and it must be cheaper than recomputing. For an 8B
  model, the break-even bandwidth is about 4\,GB/s, so even a local SSD is
  usually fast enough.
\item \textbf{Cache-aware routing}: send the user back to the replica that
  already holds the history. \emph{Limit:} replicas holding popular
  prefixes get overloaded, so routing must balance cache hits against
  load.
\end{itemize}

\begin{rubric}
\item 0.25 each for two correct techniques, 0.25 for a correct limit on at
  least one
\end{rubric}

% =====================================================================
\newpage
\section*{Set C: The ideas behind your A6 code}

\subsection*{5(a) Left padding}

\textbf{(i)} With left padding, every row's \textbf{last token is at the
same slot} (slot 4), so at each decode step every row's new token goes to
the same next slot (5, 6, \ldots): the cache grows by one column for the
whole batch. With right padding, X's next token would go to slot 3 while
Y's goes to slot 5. The rows would write to different slots, and X's new
tokens would land in the middle of its padding.

\textbf{(ii)} \textbf{Position 0}, because GPT-2 was trained with every
text starting at position 0, and X must look the same as when it runs
alone. So X's real tokens at slots 2, 3, 4 get positions 0, 1, 2
(position = slot $-$ pad). X's tokens must never attend to \textbf{slots 0
and 1}, the padding: they hold no real token. The pad mask hides them,
and the causal mask hides the future slots.

\textbf{(iii)} GPT-2 adds a learned \textbf{position embedding} to each
token. With slot numbers, X's text ``starts'' at position 2, so every one
of X's tokens gets a different embedding than when X runs alone, and the
attention results change, so the generated tokens change. Nothing crashes
because positions 2--4 are valid. Y has no padding, so its slot numbers
and positions are the same, and Y is not affected.

\begin{rubric}
\item 0.25: left padding lines up every row's newest token, so all rows
  write to the same slot
\item 0.25: X starts at position 0, and the pad mask hides slots 0--1
\item 0.25: the shifted position embeddings explain why X is wrong and Y
  is right
\end{rubric}

\subsection*{5(b) Blocks that are already promised}

\textbf{(i)} A request's full length in blocks is its tokens divided by
16, rounded up:
\begin{align*}
\text{free now} &= 20 - 4 - 2 = 14 \\
\text{still promised to R1} &= \lceil 100/16 \rceil - 4 = 7 - 4 = 3 \\
\text{still promised to R2} &= \lceil 64/16 \rceil - 2 = 4 - 2 = 2 \\
\text{A needs } \lceil 120/16 \rceil = 8 &\;\le\; 14 - 5 = 9
  \;\Rightarrow\; \textbf{admit A} \\
\text{B needs } \lceil 60/16 \rceil = 4 &\;>\; 14 - 5 - 8 = 1
  \;\Rightarrow\; \textbf{B waits}
\end{align*}
Once A is admitted, all 8 of its blocks are promised too, even though it
holds none yet.

\textbf{(ii)} Under the simpler rule, A ($8 \le 14$) and B ($4 \le 14$)
are both admitted, but at full length the four requests need $7 + 4 + 8
+ 4 = 23$ blocks, more than the 20 in the pool. It does not fail at
admission because blocks are handed out \textbf{as sequences grow}. The 14
``free'' blocks include 5 that R1 and R2 will claim later, so ``free right
now'' overstates the room. The pool runs dry only after the sequences
have grown, many decode steps later.

Safe admission \textbf{promises} that every admitted request can finish
without being evicted, so it has no preemption plan: running out means
the promise was broken, which is a bug, and the engine stops with an
error. Optimistic admission overcommits on purpose: it expects to run
out, so it \textbf{preempts} a request, frees its blocks, and recomputes
it later with one prefill.

\begin{rubric}
\item 0.25: free 14, 5 still promised, A admitted ($8 \le 9$), B waits
  ($4 > 1$, counting A's 8)
\item 0.25: blocks are handed out as sequences grow, so ``free now''
  includes blocks already promised; the shortage shows up later
\item 0.25: safe mode promises no preemption, so running out is a bug;
  optimistic mode expects it and preempts
\end{rubric}

\textbf{Common mistake.} Forgetting A's 8 blocks when B is checked:
A holds none yet, but they are already promised.

\end{document}
